\chapter{Improving the score on a fixed model, and gated scene adaptation}\label{ch:scoringgains}
\begingroup\sloppy % long provenance tags need loose lines rather than overfull ones

\begin{onething}
On 2026-10-19 and 20 four ways to find more targets without a better model were tried on FT04's fixed reconstruction -
blurring the residual (a spatial matched filter), fusing residual Mahalanobis with PCA-RX by rank, whitening each
background type separately, and briefly retraining on each scene with the likely anomalies gated out of the loss - and
every one lost real objects to plain residual Mahalanobis. A pixel measure counts a big object's pixels many times and
a one-pixel target once, so pass criteria must count objects caught; and a pretrained model has no early window for
GT-HAD's early stopping to find.
\end{onething}

\section{The question: change the score, not the model}\label{sec:scoringgains:question}

\Cref{ch:ftruns} closed on a hard fact. v0.2.0, FT03 and FT04 all tie on detection once you score them with residual
Mahalanobis (RM, \cref{ch:residualmaha}), and \cref{sec:ftruns:absorb} showed why: if a better model shrinks every
residual by the same factor, the whitening shrinks with it and every score stays put. The experiment record puts it in
three words: RM ``absorbs model gains''. So the question on 2026-10-19 was the other half: \textbf{keep the
reconstruction fixed and change only the score.} Can a better score find more?

These are evaluation scorers in \code{app/evaluation/scorers.py}, run by \file{configs/eval/scoring_gains.json} on one
model (FT04 epoch 10) and two datasets. Nothing in the product's \code{anomaly_scoring} changed, and none of them was
adopted. The baseline throughout is FT04 with plain RM and \emph{no} cloud exclusion - pAUC 0.162, not the 0.170 of
Indradhanu V2 with its cloud mask - so every row is on the same footing.

\begin{analogy}
A school wants better exam results without changing the teaching, so it tries changing the marking. Rule one: give credit
for answers that are nearly right. The class average rises - but the students who got one hard question exactly right and
nothing else now look like everyone else. Rule two: average two examiners' rankings. Rule three: grade each class against
its own class average. Rule three makes the mock exam look great, because the mock spread its hard questions evenly over
every class, while the real exam's hard cases all sat in one small class - who, graded against themselves, now look
ordinary. Before trusting any new marking, count how many individual students it identifies correctly, not only the
average mark. Each rule is one section of this chapter.
\end{analogy}

\paragraph{The pass criteria, fixed before the run.} A scorer counted as a gain over RM on FT04 if:
\begin{enumerate}
\item real PRISMA bench: pAUC up to 1\% false alarms rose by at least 0.005 \emph{and} TPR at 0.1\% false alarms did not
  fall;
\item bench-v2.1: objects caught at the top 1\% fell by no more than 0.5 points.
\end{enumerate}
Criterion 2 is only a guard. bench-v2.1's implants are blurred with $\sigma = 1.4$ by construction (\cref{ch:benchmarks}),
so a blur of 1.4 would partly be measuring its own reflection there. The real bench is the judge.

\begin{table}[htbp]
\centering
\caption{The scorers on FT04's fixed reconstruction. Real bench: 4 PRISMA scenes, 909 objects; bench-v2.1: 18 hosts.}
\label{tab:scoringgains:results}
\small\setlength{\tabcolsep}{4pt}
\begin{tabular}{lccccc}
\toprule
 & \multicolumn{3}{c}{real PRISMA bench} & bench-v2.1 & \\
\cmidrule(lr){2-4}\cmidrule(lr){5-5}
scorer & pAUC & TPR @ 0.1\% & caught 1\% / 0.1\% & caught 1\% & criteria \\
\midrule
RM (baseline) & 0.162 & 0.061 & \textbf{45.7\% / 14.9\%} & 35.1\% & -- \\
A1 RM blur 1.0 & 0.177 & 0.064 & 40.7\% / 12.0\% & 33.4\% & 1 met, 2 failed \\
A1 RM blur 1.4 & 0.177 & 0.069 & 36.3\% / 10.2\% & 29.4\% & 1 met, 2 failed \\
A1+ RM + RM blur 1.4 & 0.180 & 0.064 & 43.0\% / 12.8\% & 35.0\% & both met \\
A2 rank mean & 0.162 & 0.057 & 43.8\% / 14.0\% & 34.9\% & failed \\
A2 rank max & 0.138 & 0.046 & 42.7\% / 11.4\% & 34.8\% & failed \\
A3 mixture $k = 6$ & 0.111 & 0.045 & 39.4\% / 11.4\% & 41.2\% & failed \\
A3 mixture $k = 10$ & 0.096 & 0.034 & 34.4\% / 8.1\% & 41.9\% & failed \\
PCA-RX (no model) & 0.128 & 0.029 & 39.7\% / 8.1\% & 33.2\% & -- \\
\bottomrule
\end{tabular}

\smallskip
{\footnotesize Source: \tagM{research/experiments/scoring-gains/README.md}; full tables in \file{research/experiments/scoring-gains/report.md}.}
\end{table}

Read the bold column first. \textbf{No scorer catches as many real objects as plain RM, at either budget.}
\Cref{fig:scoringgains:scatter} draws the puzzle to scale: the A1 points move right on pixel pAUC and down on objects
caught.

\begin{figure}[htbp]
\centering
\begin{tikzpicture}
\begin{axis}[width=0.85\linewidth, height=6.8cm, xlabel={real-bench pAUC to 1\% (pixels)},
  ylabel={real objects caught at top 1\% (\%)}, xmin=0.09, xmax=0.19, ymin=33, ymax=48,
  scaled x ticks=false, xticklabel style={/pgf/number format/fixed, /pgf/number format/precision=2},
  grid=major, grid style={gray!20}, every node near coord/.append style={font=\scriptsize},
  nodes near coords, point meta=explicit symbolic]
\addplot[only marks, mark=*, mark size=2.5pt, color=spat, nodes near coords align={anchor=south west}]
  coordinates {(0.162,45.7)[RM] (0.128,39.7)[PCA-RX]};
\addplot[only marks, mark=square*, mark size=2.5pt, color=loss, nodes near coords align={anchor=north west}]
  coordinates {(0.177,40.7)[blur 1.0] (0.177,36.3)[blur 1.4] (0.180,43.0)[A1+]};
\addplot[only marks, mark=triangle*, mark size=3pt, color=attn, nodes near coords align={anchor=south east}]
  coordinates {(0.162,43.8)[mean] (0.138,42.7)[max]};
\addplot[only marks, mark=diamond*, mark size=3pt, color=spec, nodes near coords align={anchor=west}]
  coordinates {(0.111,39.4)[$k$=6] (0.096,34.4)[$k$=10]};
\draw[dashed, maskgrey] (axis cs:0.162,33) -- (axis cs:0.162,48);
\draw[dashed, maskgrey] (axis cs:0.09,45.7) -- (axis cs:0.19,45.7);
\end{axis}
\end{tikzpicture}
\caption{Pixel pAUC against objects caught, real PRISMA bench, drawn to scale from \cref{tab:scoringgains:results}. The
dashed lines cross at plain RM. Nothing lies above the horizontal line: the A1 blurs (amber) move right and down - more
pixel pAUC, fewer objects.}
\label{fig:scoringgains:scatter}
\end{figure}

\section{A1: the spatial matched filter}\label{sec:scoringgains:a1}

\subsection{The matched filter, derived}

\index{matched filter}
A \emph{matched filter} is a weighting shaped like the signal you are looking for. Set it up in general. A measurement
$\vect{x} \in \R^n$ is a known signal shape $\vect{t}$ with unknown amplitude $a$, plus noise $\vect{n}$ with mean
$\vect{0}$ and covariance $\Cov$:
\[
\vect{x} = a\,\vect{t} + \vect{n}.
\]
A linear filter $\vect{w}$ turns it into one number, $y = \vect{w}\T\vect{x} = a\,\vect{w}\T\vect{t} + \vect{w}\T\vect{n}$.
The signal part is $a\,\vect{w}\T\vect{t}$. The noise part has mean 0 and variance
\[
\operatorname{Var}[\vect{w}\T\vect{n}] = \E[(\vect{w}\T\vect{n})^2] = \E[\vect{w}\T\vect{n}\vect{n}\T\vect{w}] =
\vect{w}\T\E[\vect{n}\vect{n}\T]\vect{w} = \vect{w}\T\Cov\vect{w}.
\]
(The second step writes a square as a number times its own transpose; the third moves the constant $\vect{w}$ out of
the expectation.) The \emph{signal-to-noise ratio} per unit amplitude is
\begin{equation}
\operatorname{SNR}(\vect{w}) = \frac{\vect{w}\T\vect{t}}{\sqrt{\vect{w}\T\Cov\vect{w}}}. \label{eq:scoringgains:snr}
\end{equation}
Which $\vect{w}$ maximises it?

\paragraph{Step 1: the scale does not matter.} Replace $\vect{w}$ by $c\,\vect{w}$ with $c > 0$: the numerator gains a
factor $c$ and the square root of the denominator gains $\sqrt{c^2} = c$. They cancel. So fix the scale by demanding
$\vect{w}\T\vect{t} = 1$, and maximising the SNR becomes \emph{minimising the noise variance} $\vect{w}\T\Cov\vect{w}$
subject to $\vect{w}\T\vect{t} = 1$.

\paragraph{Step 2: a Lagrange multiplier.} \index{Lagrange multiplier}
To minimise a function subject to a constraint, minimise the \emph{Lagrangian}, the function plus a multiplier $\lambda$
times the constraint written as ``$= 0$'':
\begin{align}
\mathcal{L}(\vect{w}, \lambda) &= \vect{w}\T\Cov\vect{w} - \lambda\,(\vect{w}\T\vect{t} - 1).
\intertext{Differentiate with respect to $\vect{w}$. For a symmetric $\Cov$, the gradient of $\vect{w}\T\Cov\vect{w}$ is
$2\Cov\vect{w}$ (in one dimension, $\frac{d}{dw}\sigma^2 w^2 = 2\sigma^2 w$), and the gradient of $\vect{w}\T\vect{t}$ is
$\vect{t}$. Set the gradient to zero:}
2\Cov\vect{w} - \lambda\vect{t} &= \vect{0}.
\intertext{Multiply both sides on the left by $\tfrac12\Cov\inv$:}
\vect{w} &= \tfrac{\lambda}{2}\,\Cov\inv\vect{t}.
\intertext{Put this into the constraint $\vect{t}\T\vect{w} = 1$ to find $\lambda$: $\tfrac{\lambda}{2}\,\vect{t}\T\Cov\inv
\vect{t} = 1$, so $\tfrac{\lambda}{2} = 1 / (\vect{t}\T\Cov\inv\vect{t})$. Therefore}
\vect{w}^\star &= \frac{\Cov\inv\vect{t}}{\vect{t}\T\Cov\inv\vect{t}}. \label{eq:scoringgains:mf}
\end{align}
Its SNR follows by substitution. The numerator is $\vect{w}^{\star\T}\vect{t} = 1$ by construction. The denominator:
$\vect{w}^{\star\T}\Cov\vect{w}^\star = \frac{\vect{t}\T\Cov\inv\Cov\Cov\inv\vect{t}}{(\vect{t}\T\Cov\inv\vect{t})^2} =
\frac{1}{\vect{t}\T\Cov\inv\vect{t}}$. So
\begin{equation}
\operatorname{SNR}(\vect{w}^\star) = \sqrt{\vect{t}\T\Cov\inv\vect{t}}, \label{eq:scoringgains:mfsnr}
\end{equation}
the Mahalanobis length of the signal shape. No linear filter does better: that is what ``matched'' means.

\begin{worked}[label={we:scoringgains:mf2}]{A two-band matched filter}
Two bands with correlated noise, $\Cov = \begin{psmallmatrix} 2 & 1 \\ 1 & 2 \end{psmallmatrix}$, and a signal only in band
1, $\vect{t} = (1, 0)$. Do it on paper.

\emph{Invert $\Cov$.} For a $2 \times 2$ matrix $\begin{psmallmatrix} a & b \\ c & d\end{psmallmatrix}$ the inverse is
$\frac{1}{ad - bc}\begin{psmallmatrix} d & -b \\ -c & a\end{psmallmatrix}$. Here $ad - bc = 2 \times 2 - 1 \times 1 = 3$:
\[
\Cov\inv = \tfrac13 \begin{pmatrix} 2 & -1 \\ -1 & 2 \end{pmatrix}.
\]
\emph{Apply to $\vect{t}$.} $\Cov\inv\vect{t} = \tfrac13 (2, -1) = (0.667, -0.333)$, and
$\vect{t}\T\Cov\inv\vect{t} = 1 \times 0.667 + 0 \times (-0.333) = 0.667$.

\emph{The filter.} $\vect{w}^\star = (0.667, -0.333) / 0.667 = (1, -0.5)$. It takes band 1 and \emph{subtracts half of band
2}: band 2 holds no signal, but its noise is correlated with band 1's, so subtracting part of it cancels noise.

\emph{Check the SNR.} $\Cov\vect{w}^\star = (2 \times 1 + 1 \times (-0.5),\ 1 \times 1 + 2 \times (-0.5)) = (1.5, 0)$;
$\vect{w}^{\star\T}\Cov\vect{w}^\star = 1 \times 1.5 + (-0.5) \times 0 = 1.5$; SNR $= 1 / \sqrt{1.5} = 0.816$, which equals
$\sqrt{0.667} = 0.816$, as \eqref{eq:scoringgains:mfsnr} says. The naive filter $\vect{w} = \vect{t} = (1, 0)$ gives
$\vect{w}\T\Cov\vect{w} = 2$ and SNR $= 1/\sqrt{2} = 0.707$. Matching gains about 15\%. (Script check\tagC.)
\end{worked}

\paragraph{The spatial version.} Now let $\vect{x}$ be one band's residual over a small neighbourhood of pixels rather than
several bands at one pixel. If the residual noise is independent from pixel to pixel with equal variance, $\Cov = \sigma^2
\mat{I}$ and \eqref{eq:scoringgains:mf} gives $\vect{w}^\star \propto \vect{t}$: \textbf{weight the neighbours by the
target's own footprint}. A target that always arrives smeared by the sensor's point-spread function (PSF, $\sigma \approx
1.4$\,px on both benches; \cref{ch:benchmarks}) is matched by a Gaussian blur of the same width. That is A1: blur the
residual (\code{spatial_sigma} 1.0 or 1.4) before whitening.

\begin{worked}[label={we:scoringgains:tap}]{The promise and its price, on a 3-tap kernel}
One band, already in units of the background's spread (noise variance 1, mean 0). Use the kernel $[0.25, 0.5, 0.25]$ as a
stand-in for the Gaussian.

\emph{Noise after the blur.} The variance of a weighted sum of independent values is the sum of squared weights times each
variance:
\begin{align*}
0.25^2 + 0.5^2 + 0.25^2 &= 0.0625 + 0.25 + 0.0625 = 0.375\\
\text{spread} &= \sqrt{0.375} = 0.612
\end{align*}
\emph{A footprint-shaped target} $[1, 2, 1]$:
\begin{align*}
\text{blurred centre} &= 0.25 \times 1 + 0.5 \times 2 + 0.25 \times 1 = 0.25 + 1 + 0.25 = 1.5\\
\text{SNR raw} &= 2 / 1 = 2.00, \qquad \text{SNR blurred} = 1.5 / 0.612 = 2.45
\end{align*}
The kernel is proportional to $[1, 2, 1]$, so it \emph{is} this target's matched filter, and 2.45 equals
$\sqrt{\vect{t}\T\vect{t}} = \sqrt{1 + 4 + 1} = \sqrt{6} = 2.449$ from \eqref{eq:scoringgains:mfsnr}.

\emph{A one-pixel target} $[0, 2, 0]$:
\begin{align*}
\text{blurred centre} &= 0.25 \times 0 + 0.5 \times 2 + 0.25 \times 0 = 1.0\\
\text{SNR blurred} &= 1.0 / 0.612 = 1.63 \qquad\text{(raw was 2.00)}
\end{align*}
The same filter that helps a footprint-shaped target hurts a sharper one: a matched filter is matched to \emph{one} shape.
(Script: 2.449 and 1.633\tagC.)
\end{worked}

In the real code (\code{blurred_mahalanobis}) the whitening is refitted on the blurred residual - the ``divide by the new
spread'' step above. And \code{blur} uses \emph{normalised convolution}, so nodata does not bleed into the scene's edge:

\begin{lstlisting}
# app/evaluation/scorers.py, blur
k = valid.astype(np.float32)                               # (H, W) 1 valid, 0 nodata
w = np.maximum(gaussian_filter(k, sigma), 1e-6)            # blurred validity
out[i] = gaussian_filter(x[i] * k, sigma) / w * k          # per band (H, W)
\end{lstlisting}

\index{normalised convolution}
Blur the data with nodata set to 0, then divide by the blurred validity. At a pixel next to nodata, the blurred data
counts fewer real neighbours, and dividing by the blurred validity renormalises the weights over the valid ones only. With
the 3-tap kernel, values $(2, 4, \text{nodata})$ and the centre pixel: data blur $0.25 \times 2 + 0.5 \times 4 + 0.25
\times 0 = 2.5$; validity blur $0.25 + 0.5 + 0 = 0.75$; result $2.5 / 0.75 = 3.33$, the average of 2 and 4 with weights
$0.25/0.75$ and $0.5/0.75$. Without the division you would get 2.5, dragged towards the zeros.

\subsection{What a pixel measure rewards}\label{sec:scoringgains:pixels}

\index{pixel measure}\index{object measure}
pAUC and TPR are \emph{pixel measures}: they count target pixels above a threshold. ``Objects caught'' is an \emph{object
measure}: an object counts once if any of its pixels is near a top-scoring pixel (\cref{ch:measures}). A filter that fills
in more pixels of an object already found raises the first and not the second. You need a pen for the next example; it
takes about five minutes.

\begin{worked}[label={we:scoringgains:strip}]{A strip of 12 pixels, plain and blurred}
One band, background mean 0, score $=$ residual squared (a one-band global RM). Pixels 5--8 are a big object, pixel 11 is
a one-pixel target, pixels 2 and 3 are noise that flips sign from one pixel to the next.
\begin{center}
\small\setlength{\tabcolsep}{3.5pt}
\begin{tabular}{lcccccccccccc}
\toprule
pixel & 1 & 2 & 3 & 4 & 5 & 6 & 7 & 8 & 9 & 10 & 11 & 12 \\
\midrule
residual & 0 & 4.5 & $-4.2$ & 0 & 3 & 4 & 5 & 3 & 0 & 0 & 4.8 & 0 \\
score $r^2$ & 0 & 20.25 & 17.64 & 0 & 9 & 16 & 25 & 9 & 0 & 0 & 23.04 & 0 \\
\bottomrule
\end{tabular}
\end{center}
The analyst's budget is the top 4 pixels. Plain RM's top 4: pixel 7 (25), 11 (23.04), 2 (20.25), 3 (17.64).

Blur each pixel, $b_i = 0.25\,r_{i-1} + 0.5\,r_i + 0.25\,r_{i+1}$, zero outside the strip:
\begin{align*}
b_2 &= 0.25 \times 0 + 0.5 \times 4.5 + 0.25 \times (-4.2) = 0 + 2.25 - 1.05 = 1.2 &&\to 1.44\\
b_3 &= 0.25 \times 4.5 + 0.5 \times (-4.2) + 0.25 \times 0 = 1.125 - 2.1 + 0 = -0.975 &&\to 0.95\\
b_5 &= 0.25 \times 0 + 0.5 \times 3 + 0.25 \times 4 = 0 + 1.5 + 1 = 2.5 &&\to 6.25\\
b_6 &= 0.25 \times 3 + 0.5 \times 4 + 0.25 \times 5 = 0.75 + 2 + 1.25 = 4.0 &&\to 16\\
b_7 &= 0.25 \times 4 + 0.5 \times 5 + 0.25 \times 3 = 1 + 2.5 + 0.75 = 4.25 &&\to 18.06\\
b_8 &= 0.25 \times 5 + 0.5 \times 3 + 0.25 \times 0 = 1.25 + 1.5 + 0 = 2.75 &&\to 7.56\\
b_{11} &= 0.25 \times 0 + 0.5 \times 4.8 + 0.25 \times 0 = 2.4 &&\to 5.76
\end{align*}
(Pixels 1, 4, 9, 10 and 12 come out at 1.27, 0.09, 0.56, 1.44 and 1.44 - all small.) Blurred top 4: pixel 7 (18.06), 6
(16), 8 (7.56), 5 (6.25). The target, at 5.76, is fifth.
\begin{center}
\small
\begin{tabular}{lccc}
\toprule
 & object pixels in top 4 (of 5) & false alarms in top 4 & objects caught (of 2) \\
\midrule
plain & 2 (pixels 7, 11) & 2 (pixels 2, 3) & 2 \\
blurred & 4 (pixels 5--8) & 0 & 1 \\
\bottomrule
\end{tabular}
\end{center}
A pixel measure such as TPR doubles, $2/5 = 0.4$ to $4/5 = 0.8$, and the false alarms vanish. Objects caught halves.
(Every value matches the script\tagC.)
\end{worked}

\begin{keypoint}
A pixel measure counts each pixel of an object, so filling in more of an object you had already found looks like progress,
and losing a one-pixel target costs one pixel's worth.
\end{keypoint}

The example simplifies three things - one band, a global mean instead of the 21/61 ring, a 3-tap kernel - but the real
bench shows the same mechanism (\cref{tab:scoringgains:size}). Blur loses objects at \emph{every} size. Turned into counts:
$0.365 \times 561 = 204.8 \approx 205$ small objects for RM against $0.333 \times 561 = 186.8 \approx 187$ blurred, 18
lost; $0.524 \times 170 = 89.1 \approx 89$ against $0.453 \times 170 = 77.0 \approx 77$ in the 3--5\,px class, 12 lost. In
the largest class the loss is one object ($0.80 \times 20 = 16$ against $0.75 \times 20 = 15$), too few to read. And 561 of
909 real objects, 62\%, are 1--2 pixels: a filter that favours footprint-sized targets works against most of the bench.

\begin{table}[htbp]
\centering
\caption{Real-bench objects caught at the top 1\%, by object size (561 / 170 / 158 / 20 objects).}
\label{tab:scoringgains:size}
\small
\begin{tabular}{lcccc}
\toprule
scorer & 1--2\,px & 3--5\,px & 6--20\,px & $>$ 20\,px \\
\midrule
RM (baseline) & \textbf{36.5\%} & \textbf{52.4\%} & \textbf{66.5\%} & 80.0\% \\
A1 RM blur 1.0 & 33.3\% & 45.3\% & 57.6\% & 75.0\% \\
A1+ RM + RM blur 1.4 & 35.1\% & 48.2\% & 60.8\% & 80.0\% \\
A3 mixture $k = 6$ & 32.6\% & 43.5\% & 54.4\% & 75.0\% \\
\bottomrule
\end{tabular}

\smallskip
{\footnotesize Source: \tagM{research/experiments/scoring-gains/README.md}.}
\end{table}

\paragraph{A1+ passed its criteria and was still not adopted.} A1+ (rank mean of RM and RM blurred at 1.4) met both
criteria: pAUC $0.180 - 0.162 = +0.018$, at least $+0.005$; TPR 0.064, not below 0.061; bench-v2.1 $35.0 - 35.1 = -0.1$
points, within 0.5. It also lost $45.7 - 43.0 = 2.7$ points of real objects at the top 1\% and $14.9 - 12.8 = 2.1$ at the
top 0.1\%. The README: ``The criteria were pixel-level and missed this; not adopted.'' A criterion that was right to fix in
advance and wrong in what it measured. From then on, every pass criterion in the record includes real objects caught.

\section{A2: rank fusion with PCA-RX}\label{sec:scoringgains:a2}

\index{rank fusion}
PCA-RX (local RX on the cube's top principal components, no model; \citet{reed1990adaptive}) wins where a target is foreign
to the \emph{scene}; RM wins where it is foreign to the \emph{model}. The hope was that they fail on different objects, so
combining them would catch the union. Their scores live on different scales, so they cannot be averaged directly.
\code{rank_map} puts them on one footing:
\begin{equation}
\operatorname{rank}(p) = \frac{\#\{\text{valid pixels with score} \le s(p)\}}{\#\{\text{valid pixels}\}} \quad
\text{(ties averaged, 0 outside valid)}, \label{eq:scoringgains:rank}
\end{equation}
the empirical cumulative distribution of the scores, evaluated at each pixel. \code{fuse} stacks two such maps and takes
the mean or the maximum.

\begin{worked}[label={we:scoringgains:rank}]{Rank mean and rank max on five pixels}
P1 is a real target. P2 is RM's false alarm (a field edge); P4 is PCA-RX's (a roof that is spectrally unusual but ordinary
to the model).
\begin{center}
\small
\begin{tabular}{lccccc}
\toprule
 & P1 & P2 & P3 & P4 & P5 \\
\midrule
RM score & 9 & 12 & 3 & 2 & 1 \\
PCA-RX score & 8 & 1 & 2 & 10 & 3 \\
RM rank & 4/5 & 5/5 & 3/5 & 2/5 & 1/5 \\
PCA-RX rank & 4/5 & 1/5 & 2/5 & 5/5 & 3/5 \\
rank mean & 0.8 & 0.6 & 0.5 & 0.7 & 0.4 \\
rank max & 0.8 & 1.0 & 0.6 & 1.0 & 0.6 \\
\bottomrule
\end{tabular}
\end{center}
Each rank counts the pixels at or below: RM's P1 (score 9) has 9, 3, 2, 1 at or below it, four of five, so $4/5$. The mean
for P1 is $(4/5 + 4/5)/2 = 0.8$; for P4, $(2/5 + 5/5)/2 = 7/10 = 0.7$.

With a budget of one pixel: RM alone picks P2, PCA-RX alone picks P4 - both miss. Rank mean picks P1: a pixel both scorers
think odd beats a pixel only one thinks odd. Rank max puts P2 and P4 first together (1.0) and P1 third. (Script
check\tagC.)
\end{worked}

Why does the maximum let false alarms through? Suppose, on background pixels, the two ranks behave like independent uniform
numbers $U, V$ on $[0, 1]$. At a threshold $1 - q$:
\begin{align*}
P(\max(U, V) > 1 - q) &= 1 - P(U \le 1 - q)\,P(V \le 1 - q) = 1 - (1 - q)^2 = 2q - q^2,\\
P\big(\tfrac{U + V}{2} > 1 - q\big) &= P(U + V > 2 - 2q) = \tfrac12 (2q)^2 = 2q^2 \quad (q \le \tfrac12).
\end{align*}
(The second is the area of the corner triangle with legs $2q$ in the unit square.) At $q = 0.01$, the maximum passes $2
\times 0.01 - 0.0001 = 0.0199$ of the background - nearly twice either scorer alone - while the mean passes $2 \times 0.0001
= 0.0002$. The maximum's top is roughly the \emph{union} of the two scorers' tops, false alarms included; the mean's top
needs agreement. (In the evaluator the fused map is re-ranked and the budget is a fixed share of it, so these numbers
describe the shape of the fused map, not the final false-alarm rate.)

The real bench gave the expected half and denied the hopeful half: rank max fell to pAUC 0.138, as predicted; rank mean tied
pAUC (0.162) and \emph{lost} objects (43.8\% against 45.7\%). Rank mean needs the scorers to fail on different objects, and
on the real bench they did not. PCA-RX alone catches 39.7\%; its wins on bench-v1 ``do not carry to real targets'', so
averaging in its ranks mostly dilutes RM.

\section{A3: a mixture of backgrounds}\label{sec:scoringgains:a3}

\subsection{Gaussian mixtures and responsibilities}

\index{Gaussian mixture}\index{responsibility}
One covariance for a whole scene mixes land, water and roofs. A \emph{Gaussian mixture} models the background instead as
$K$ Gaussian components, each a surface type: component $k$ has a weight $\pi_k$ (its share of the scene, $\sum_k \pi_k =
1$), a mean $\muv_k$ and a covariance $\Cov_k$. The density of a pixel $\vect{x}$ is
\begin{equation}
p(\vect{x}) = \sum_{k=1}^K \pi_k\,\mathcal{N}(\vect{x};\,\muv_k, \Cov_k). \label{eq:scoringgains:gmm}
\end{equation}
Which component does a pixel come from? Bayes' rule turns ``probability of the pixel given the component'' into
``probability of the component given the pixel'' - the \emph{responsibility}:
\begin{equation}
r_k(\vect{x}) = P(k \mid \vect{x}) = \frac{P(k)\,p(\vect{x} \mid k)}{p(\vect{x})} =
\frac{\pi_k\,\mathcal{N}(\vect{x};\,\muv_k, \Cov_k)}{\sum_j \pi_j\,\mathcal{N}(\vect{x};\,\muv_j, \Cov_j)}.
\label{eq:scoringgains:resp}
\end{equation}
A natural per-type anomaly score is then the \emph{responsibility-weighted} squared Mahalanobis distance,
\begin{equation}
s(\vect{x}) = \sum_k r_k(\vect{x})\,(\vect{x} - \muv_k)\T\Cov_k\inv(\vect{x} - \muv_k), \label{eq:scoringgains:rwm}
\end{equation}
each type's distance counted in proportion to how likely the pixel belongs to it.

\begin{worked}[label={we:scoringgains:gmm}]{Responsibilities and a weighted score in one band}
Two types: a field ($\pi_1 = 0.8$, mean 0, variance 1) and roofs ($\pi_2 = 0.2$, mean 5, variance 4). A pixel at $x = 3$.
The one-dimensional normal density is $\mathcal{N}(x; \mu, v) = \frac{1}{\sqrt{2\pi v}} e^{-(x - \mu)^2 / (2v)}$.
\begin{align*}
\mathcal{N}(3; 0, 1) &= \tfrac{1}{\sqrt{2\pi}}\,e^{-9/2} = 0.3989 \times 0.01111 = 0.004432\\
\mathcal{N}(3; 5, 4) &= \tfrac{1}{\sqrt{8\pi}}\,e^{-4/8} = 0.1995 \times 0.6065 = 0.1210\\
\pi_1 \mathcal{N}_1 &= 0.8 \times 0.004432 = 0.003545\\
\pi_2 \mathcal{N}_2 &= 0.2 \times 0.1210 = 0.02420\\
r_1 &= 0.003545 / (0.003545 + 0.02420) = 0.003545 / 0.02774 = 0.128\\
r_2 &= 0.02420 / 0.02774 = 0.872
\end{align*}
Squared distances: to the field $(3 - 0)^2 / 1 = 9$; to the roofs $(3 - 5)^2 / 4 = 1$. Weighted score:
\[
s = 0.128 \times 9 + 0.872 \times 1 = 1.15 + 0.87 = 2.02 .
\]
A single global Gaussian has mean $0.8 \times 0 + 0.2 \times 5 = 1$ and variance $\E[x^2] - 1^2 = 0.8 \times 1 + 0.2 \times
(4 + 25) - 1 = 0.8 + 5.8 - 1 = 5.6$, giving $(3 - 1)^2 / 5.6 = 0.71$. (Script: $r_1 = 0.128$, $s = 2.022$\tagC.)
\end{worked}

\subsection{What the code does: hard types by k-means}

\index{k-means}
\code{background_types} does not fit a soft mixture. It runs \emph{k-means} (\citet{lloyd1982least}'s algorithm) - the
hard-assignment limit of a mixture, in which every responsibility is 0 or 1: pick $k$ centres, give each pixel to its
nearest centre, move each centre to the mean of its pixels, repeat.

\begin{worked}[label={we:scoringgains:kmeans}]{k-means on six numbers}
Points $0, 1, 2, 8, 9, 10$; two centres starting at 0 and 2.
\begin{enumerate}
\item Assign. 0 is nearer 0; 1 is equally near both (distance 1) and the code's \code{argmin} takes the first, 0; 2, 8, 9,
  10 are nearer 2. Update: centre 1 $= (0 + 1)/2 = 0.5$; centre 2 $= (2 + 8 + 9 + 10)/4 = 29/4 = 7.25$.
\item Assign. 2 is now $\abs{2 - 0.5} = 1.5$ from centre 1 and $\abs{2 - 7.25} = 5.25$ from centre 2: it moves. Update:
  centre 1 $= (0 + 1 + 2)/3 = 1$; centre 2 $= (8 + 9 + 10)/3 = 9$.
\item Assign: nothing moves. Converged at 1 and 9. (Script check\tagC.)
\end{enumerate}
\end{worked}

\begin{sourcenote}[: ``mixture'' is hard k-means plus class-aware RM]
The README calls A3 a ``mixture of backgrounds'', but there is no soft Gaussian mixture in the code.
\code{background_types} (\code{app/evaluation/scorers.py}) runs Lloyd's k-means, fixed seed, 25 iterations, fitted on a
50,000-pixel sample, on the \emph{cube's} whitened top 10 principal components - so the types are surfaces, not residual
patterns - and assigns every valid pixel to its nearest centre. \code{score} then hands those types to
\code{class_aware_residual_mahalanobis} (\code{app/utils/anomaly_detection/residual_mahalanobis.py}), written for a
segmentation head's classes: each type gets its own whitening and a ring of same-type neighbours, and types under 5,000
valid pixels (\code{MIN_CLASS_PX}) join type 0. In the language of \eqref{eq:scoringgains:rwm}, every $r_k$ is 0 or 1.
\end{sourcenote}

\subsection{Why it won the implant bench and lost the real one}

A3 won bench-v2.1 by 6 points (35.1\% to 41.2\%) and cut the real bench's pAUC by a third (0.162 to 0.111). Here is why, by
hand.

\begin{worked}[label={we:scoringgains:types}]{Targets that become their type's normal}
One band, every type's mean 0, score $=$ residual squared $/$ variance. An illustration\tagA:
\begin{itemize}
\item field type (16 pixels): seven at $+1$, seven at $-1$, one edge at $+2.5$, one at $-2.5$;
\item built-up type (6 pixels): $-2, -2, +2, +2$ (ordinary roofs), and two real targets at $-6$ and $+6$.
\end{itemize}
\emph{One global variance, all 22 pixels.}
\begin{align*}
\text{sum of squares} &= 14 \times 1 + 2 \times 6.25 + 4 \times 4 + 2 \times 36 = 14 + 12.5 + 16 + 72 = 114.5\\
\text{variance} &= 114.5 / 22 = 5.20\\
\text{target score} &= 36 / 5.20 = 6.92, \qquad \text{edge score} = 6.25 / 5.20 = 1.20
\end{align*}
The targets top the ranking.

\emph{Each type on its own.}
\begin{align*}
\text{field:}\quad & (14 + 12.5)/16 = 26.5 / 16 = 1.656; \quad \text{edge score} = 6.25 / 1.656 = 3.77\\
\text{built-up:}\quad & (16 + 72)/6 = 88 / 6 = 14.67; \quad \text{target score} = 36 / 14.67 = 2.45
\end{align*}
Now the two field edges (3.77) outrank both targets (2.45). The targets are two of the built-up type's six pixels, so they
\emph{are} a large share of what that type's ``normal'' is fitted on. Graded against themselves, they look ordinary.
\end{worked}

That is what the README found on the real scenes: 60--70\% of target pixels sit in two or three small background types of
about 5\% of the scene each - the built-up surfaces where targets live. In set4, one type holds 60\% of the target pixels on
5.2\% of the scene\tagM{research/experiments/scoring-gains/README.md}.

And the implants? bench-v2.1 spreads them over every surface. Put an implant with residual 4 into the field type above:
globally $16 / 5.20 = 3.07$, per type $16 / 1.656 = 9.66$. A tighter normal makes an implant on an ordinary surface stand out
more - so per-type whitening only sharpens the implants.

\begin{keypoint}
bench-v2.1 has a second bias (after bench-v1's cube-level oddity, \cref{ch:benchmarks}): implant placement is not
conditioned on where real targets occur. Judge scorers on the real bench and use bench-v2.1 to catch regressions.
\end{keypoint}

\file{docs/09-known-issues.md} records the bias and says the fix (placement on target-bearing surfaces) is not built.

\section{Expectations, a cache bug, and the verdict}\label{sec:scoringgains:verdict}

The README wrote its expectations down before the run. Scored afterwards:
\begin{center}
\small
\begin{tabular}{lp{5.2cm}p{5.6cm}}
\toprule
idea & expected & what happened \\
\midrule
A1 & a small gain on the real bench & right on pixels, wrong on objects \\
A2 & rank mean a small gain, rank max a loss & rank max right; rank mean tied pAUC and lost objects \\
A3 & a gain on multi-surface scenes (set2 / set4) & wrong in direction \\
\bottomrule
\end{tabular}
\end{center}

\paragraph{A float16 overflow.} \index{float16}
The run also exposed a bug in the evaluator's cache (score maps are cached, \cref{ch:runeval}). Maps were stored as
float16, a 16-bit floating-point format with 1 sign bit, 5 exponent bits and 10 fraction bits. Its largest finite value
has the largest normal exponent, $2^{15}$, and all ten fraction bits set, a significand of $1 + (1 - 2^{-10}) = 2 -
2^{-10}$:
\[
(2 - 2^{-10}) \times 2^{15} = 2^{16} - 2^{5} = 65{,}536 - 32 = 65{,}504 .
\]
Blurred and mixture squared distances went past it, so a reloaded map had \texttt{inf} at its top - exactly the pixels that
matter. \code{runner.py} now keeps a map float32 if any value is non-finite or at least $6 \times 10^4$ in size; the 13
overflowed cache files were deleted. The README notes that this run's metrics were computed before the cast and are
unaffected.

\paragraph{Verdict: RM stays the score.} The lesson carried forward in writing: the next attempt's pass criteria include
real objects caught at the top 1\% and 0.1\%. The next attempt was on the model side - an uncertainty head, FT06, which
failed too (\cref{ch:uncertainty}): it learned to expect large errors on the targets themselves, so dividing by that
expectation hid them (pAUC 0.170 to 0.125).

\section{Gated scene adaptation}\label{sec:scoringgains:adapt}

\subsection{What GT-HAD's gating and early stopping are for}

Every failure so far had one shape: \emph{every attempt to learn ``normal'' from a scene's own context also learned that
scene's targets} - the mixture made targets their type's normal, and FT06's head learned to expect them. One method in the
record seemed to have an answer. GT-HAD \citep{lian2025gated}, a published hyperspectral anomaly detector, was reproduced on
the four PRISMA scenes (\file{research/experiments/gt-had-prisma/summary.md}). It is a reconstruction detector with no
pretraining: it trains from scratch on the very scene it scores, cut into $9 \times 9$ pixel blocks. Train long enough and
it learns to reconstruct everything, targets included. It has two defences:

\begin{itemize}
\item \textbf{Gating.} \index{gating} A periodic search (the paper calls it CMM) decides which blocks look like background;
  those are routed to a background branch that learns from them, so its idea of normal is built without the likely
  anomalies.
\item \textbf{Early stopping.} \index{early stopping} Stop training before the network has had time to learn the rare
  things: a network learns the common background before the rare targets.
\end{itemize}

The second defence is visible in the reproduction:
\begin{center}
\small
\begin{tabular}{lccc}
\toprule
 & 5 epochs & 25 epochs & 150 epochs (paper's setting) \\
\midrule
mean valid-pixel AUC & \textbf{0.873} & 0.780 & 0.843 \\
mean PD at $10^{-3}$ false alarms & \textbf{5.5\%} & -- & 2.0\% \\
\bottomrule
\end{tabular}

\smallskip
{\footnotesize Source: \tagM{research/experiments/gt-had-prisma/summary.md}.}
\end{center}
(PD, probability of detection, is the same idea as TPR.) The summary is careful: ``choosing the epoch count from these GT
AUCs would leak labels, so 5 epochs is a diagnostic, not a recommendation.'' After the GT-HAD discussion the user's comment -
``Gating is something we can use'' - started the experiment of 2026-10-20.

\subsection{The transplant}

GT-HAD starts from nothing; this version starts from FT04 and only adapts it. Per scene, starting from FT04's weights each
time (\file{research/experiments/scene-adaptation/README.md}):
\begin{enumerate}
\item score the scene with FT04 + RM (the unadapted score);
\item gate $=$ the top 1\% of that score, dilated by 2\,px (\code{gate_mask});
\item adapt $N$ steps on the scene's own tiles, gated pixels left out of the loss (\code{adapt});
\item reconstruct, score RM with cloud excluded on \emph{unadapted} FT04's masks; restore FT04's weights.
\end{enumerate}

\code{adapt} (\code{app/foundation_models/scene_adaptation.py}) uses the fine-tuning loss - L1 plus $0.5 \times$ SAM on hidden
pixels, 65\% of tokens hidden - on batches of 8 random 128\,px tiles from the scene, with an Adam optimiser. Three details
matter:
\begin{itemize}
\item \textbf{The loss mask is hidden $\times$ valid $\times$ not-gated}: \code{learn = valid & ~gate}, then
  \code{loss_mask = hidden * l[:, None].float()}. A gated pixel can still be \emph{seen} as context; it never counts towards
  the loss.
\item \textbf{Eval mode during training.} \code{model.eval()} is called before the loop, so BatchNorm keeps the statistics it
  learned on many scenes. Gradients still flow; only those stored averages are frozen.
\item \textbf{The segmentation head is outside the loss}: parameters whose names start \texttt{seg\_head.} or
  \texttt{scale\_head.} are skipped - which is why cloud is excluded on the unadapted model's masks.
\end{itemize}
The evaluator copies the full \code{state_dict} before adapting and reloads it in a \code{finally} block, so a crash
mid-scene cannot leave the next scene starting from adapted weights.

\paragraph{How far does a 1\% gate reach?} \index{dilation}
The probe printed ``gated px 4.55\%'' on set1. Why that much, and why not more? \code{binary_dilation} with no structuring
element uses the cross: each pass adds the four direct neighbours of every gated pixel. Two passes turn an $a \times b$
rectangle of gated pixels into the rectangle, plus four side strips of width 2, plus four corner triangles:
\begin{equation}
\abs{\text{gate}} = ab + \underbrace{2 \cdot 2\,(a + b)}_{\text{side strips}} + \underbrace{4 \times 1}_{\text{corners}}
= ab + 4(a + b) + 4 . \label{eq:scoringgains:dilate}
\end{equation}
(Each corner triangle at distance up to 2 holds the single diagonal pixel at offset $(1, 1)$: $r(r-1)/2 = 1$ for $r = 2$.)

\begin{worked}[label={we:scoringgains:dilate}]{Growth of a gate under two passes}
\begin{align*}
1 \times 1:\quad & 1 + 4 \times 2 + 4 = 1 + 8 + 4 = 13 &&\text{growth } 13 / 1 = 13\\
2 \times 2:\quad & 4 + 4 \times 4 + 4 = 4 + 16 + 4 = 24 &&\text{growth } 24 / 4 = 6.0\\
3 \times 3:\quad & 9 + 4 \times 6 + 4 = 9 + 24 + 4 = 37 &&\text{growth } 37 / 9 = 4.1
\end{align*}
SciPy gives 13, 24 and 37\tagC. If set1's top 1\% were all isolated pixels, the gate would cover about 13\%. It covers
4.55\%, a growth of 4.55, between the $3 \times 3$ and $2 \times 2$ cases: the top of the score comes in clumps, so a halo
costs far less than it would for scattered single pixels. More than 95\% of the scene still trains the network.
\end{worked}

\subsection{Choosing settings without ground truth}

There are two settings: the learning rate and $N$, the number of steps. Choosing either by which detects the labelled targets
best would be cheating - a setting tuned on the answers looks better than it will be on a new scene.

\paragraph{The learning rate, from the loss alone.} \code{adapt_probe.py} ran 300 steps on set1 at three learning rates and
recorded only the adaptation loss - no labels (1.6 steps per second on the laptop GPU):
\begin{center}
\small
\begin{tabular}{lcccc}
\toprule
lr & first 50 steps & steps 100--150 & last 50 steps & relative fall \\
\midrule
$10^{-5}$ & 0.0495 & 0.0409 & 0.0389 & $0.0106 / 0.0495 = 21.4\%$ \\
$3 \times 10^{-5}$ & 0.0459 & 0.0382 & 0.0368 & $0.0091 / 0.0459 = 19.8\%$ \\
$10^{-4}$ & 0.0461 & 0.0365 & 0.0354 & $0.0107 / 0.0461 = 23.2\%$ \\
\bottomrule
\end{tabular}

\smallskip
{\footnotesize Losses: \tagM{research/experiments/scene-adaptation/README.md}; relative falls computed here\tagC.}
\end{center}
The choice was $3 \times 10^{-5}$: ``as fast as 1e-4 early, a third of the step size - the aim is to learn background without
reaching the targets.'' Its first-50 average (0.0459) matches $10^{-4}$'s (0.0461) with a third of the step. That is a
judgement, not a measurement: the loss shows how fast the background is fitted, not whether the targets are being learned,
because they are a tiny share of it (\cref{sec:scoringgains:lessons}).

\paragraph{N, by leave-one-scene-out.} \index{leave-one-scene-out}\index{cross-validation}
Four variants were run: $N = 50$, 200 and 1000 gated, and $N = 200$ with the gate at 0\% to isolate the gate's effect.
\emph{Cross-validation} estimates how a choice made on some data will do on new data by holding data out of the choice.
\emph{Leave-one-scene-out} is its version with scenes as the units: to report scene $s$, choose $N$ using only the other
scenes, then read $s$ at that $N$; repeat for every scene and average.

Why bother? Pick $N$ by the best mean over all four scenes and each reported scene has helped choose its own setting - the
reported number is the maximum of several noisy numbers, which is biased upwards. Leave-one-scene-out guarantees that no
scene's number comes from a setting tuned on that scene. It is not label-free - the other scenes' labels are used - but it is
honest about how a setting chosen on old scenes would do on a new one, which is the product's situation.

\begin{table}[htbp]
\centering
\caption{Real PRISMA bench, cloud excluded on unadapted FT04's masks.}
\label{tab:scoringgains:adapt}
\small
\begin{tabular}{lcccc}
\toprule
 & pAUC $\le$ 1\% & caught top 1\% & caught top 0.1\% & TPR @ 0.1\% \\
\midrule
FT04 (the bar) & \textbf{0.170} & 46.0\% & \textbf{15.8\%} & \textbf{0.065} \\
gated, $N = 50$ & 0.160 & 46.8\% & 15.2\% & 0.059 \\
gated, $N = 200$ & 0.159 & 46.9\% & 14.9\% & 0.059 \\
gated, $N = 1000$ & 0.161 & 46.5\% & 15.0\% & 0.057 \\
ungated, $N = 200$ & 0.159 & 46.9\% & 14.9\% & 0.059 \\
\bottomrule
\end{tabular}

\smallskip
{\footnotesize Source: \tagM{research/experiments/scene-adaptation/README.md} (from \texttt{eval\_report.md}).}
\end{table}

The pass criteria, fixed before any detection result, asked for a leave-one-scene-out pAUC gain of at least $+0.005$ over
FT04 with cloud excluded (0.170), objects caught at the top 1\% and 0.1\% not lower, and a bench-v2.1 guard. Part of the
previous section's lesson is visible: objects caught are in the criteria now.

\begin{worked}[label={we:scoringgains:loso}]{Leave-one-scene-out on the real numbers}
Per-scene pAUC, cloud excluded, from \file{eval_per_scene.csv}, rounded to three places:
\begin{center}
\small
\begin{tabular}{lcccc}
\toprule
 & set1 & set2 & set3 & set4 \\
\midrule
FT04 & 0.246 & 0.128 & 0.084 & 0.220 \\
$N = 50$ & 0.223 & 0.142 & 0.084 & 0.192 \\
$N = 200$ & 0.228 & 0.142 & 0.081 & 0.185 \\
$N = 1000$ & 0.240 & 0.139 & 0.083 & 0.181 \\
\bottomrule
\end{tabular}
\end{center}
The baseline: FT04 mean $= (0.246 + 0.128 + 0.084 + 0.220)/4 = 0.678/4 = 0.1695$. You need a pen; it takes about ten minutes.

\emph{Hold out set1} (mean over set2, set3, set4):
\begin{align*}
N = 50:&\quad (0.142 + 0.084 + 0.192)/3 = 0.418/3 = 0.1393\\
N = 200:&\quad (0.142 + 0.081 + 0.185)/3 = 0.408/3 = 0.1360\\
N = 1000:&\quad (0.139 + 0.083 + 0.181)/3 = 0.403/3 = 0.1343
\end{align*}
Best $N = 50$; set1 reports its $N = 50$ value, \textbf{0.223}.

\emph{Hold out set2} (set1, set3, set4):
\begin{align*}
N = 50:&\quad (0.223 + 0.084 + 0.192)/3 = 0.499/3 = 0.1663\\
N = 200:&\quad (0.228 + 0.081 + 0.185)/3 = 0.494/3 = 0.1647\\
N = 1000:&\quad (0.240 + 0.083 + 0.181)/3 = 0.504/3 = 0.1680
\end{align*}
Best $N = 1000$; set2 reports \textbf{0.139}.

\emph{Hold out set3} (set1, set2, set4):
\begin{align*}
N = 50:&\quad (0.223 + 0.142 + 0.192)/3 = 0.557/3 = 0.1857\\
N = 200:&\quad (0.228 + 0.142 + 0.185)/3 = 0.555/3 = 0.1850\\
N = 1000:&\quad (0.240 + 0.139 + 0.181)/3 = 0.560/3 = 0.1867
\end{align*}
Best $N = 1000$; set3 reports \textbf{0.083}.

\emph{Hold out set4} (set1, set2, set3):
\begin{align*}
N = 50:&\quad (0.223 + 0.142 + 0.084)/3 = 0.449/3 = 0.1497\\
N = 200:&\quad (0.228 + 0.142 + 0.081)/3 = 0.451/3 = 0.1503\\
N = 1000:&\quad (0.240 + 0.139 + 0.083)/3 = 0.462/3 = 0.1540
\end{align*}
Best $N = 1000$; set4 reports \textbf{0.181}.

\emph{The headline.}
\begin{align*}
\text{leave-one-scene-out} &= (0.223 + 0.139 + 0.083 + 0.181)/4 = 0.626/4 = 0.1565\\
\text{gain over FT04} &= 0.1565 - 0.1695 = -0.013
\end{align*}
The criterion asked for $+0.005$. On the unrounded CSV the same picks give 0.1567\tagC; the README records 0.157.
\end{worked}

Notice two things. Set1 picked a different $N$ from the others: three scenes are a small jury, and the choice is not stable.
And the honest number, 0.1565, is a little worse than the best single row's mean (0.161 for $N = 1000$). \textbf{That gap is
the size of the flattery you would have reported by choosing $N$ on all four scenes.}

\begin{figure}[htbp]
\centering
\begin{tikzpicture}
\begin{axis}[ybar, bar width=6pt, width=0.9\linewidth, height=6cm, ymin=0, ymax=0.27,
  ylabel={pAUC to 1\% (cloud excluded)}, symbolic x coords={set1,set2,set3,set4}, xtick=data,
  enlarge x limits=0.18, legend style={at={(0.5,1.02)},anchor=south,legend columns=4,font=\small},
  every axis plot/.append style={draw=black!50}]
\addplot[fill=spat!60] coordinates {(set1,0.246) (set2,0.128) (set3,0.084) (set4,0.220)};
\addplot[fill=loss!40] coordinates {(set1,0.223) (set2,0.142) (set3,0.084) (set4,0.192)};
\addplot[fill=loss!70] coordinates {(set1,0.228) (set2,0.142) (set3,0.081) (set4,0.185)};
\addplot[fill=loss] coordinates {(set1,0.240) (set2,0.139) (set3,0.083) (set4,0.181)};
\legend{FT04, $N = 50$, $N = 200$, $N = 1000$}
\end{axis}
\end{tikzpicture}
\caption{Per-scene pAUC of FT04 and its gated adaptations, drawn to scale. Set2 gains at every $N$; set1 and set4 lose at
every $N$; nothing peaks at small $N$.}
\label{fig:scoringgains:perscene}
\end{figure}

\subsection{Results, and two lessons}\label{sec:scoringgains:lessons}

Read \cref{tab:scoringgains:adapt} in three passes. Every $N$ is below FT04 on pAUC (about $-0.01$) and on objects caught at
the top 0.1\% ($-0.6$ to $-0.9$ points): criteria 1 and 2 fail. Objects caught at the top 1\% rose by 0.5 to 0.9 points;
with the top 0.1\% falling, the README reads this as ``a shift of hits from the very top into the top 1\%, not more finds.''
And gated and ungated $N = 200$ are the same row, to three decimals, on every measure. Both written expectations - a gain
peaking at small $N$, and ungated worse than gated - were wrong. bench-v2.1 was not run; with criteria 1 and 2 failed it could
not change the verdict. One scene, set2, gains at every $N$ (0.128 to 0.139--0.142); ``four scenes cannot say whether that is
signal.''

\paragraph{Lesson 1: the gate changed nothing, because adaptation never learned the targets.}
If gated and ungated agree to three decimals, the ungated network was not learning the targets either, so the gate had
nothing to prevent. Why not? Look at their weight in the loss. GT-HAD's summary gives the scale: scenes of about $1200
\times 1200$ pixels with 600 to 1200 anomaly pixels.
\begin{align*}
\text{scene pixels} &= 1{,}200 \times 1{,}200 = 1{,}440{,}000\\
\text{anomaly pixels} &\approx 1{,}000\\
\text{share} &= 1{,}000 / 1{,}440{,}000 = 0.0007 = 0.07\%
\end{align*}
In a loss that is a mean over pixels, 0.07\% of the pixels contribute about 0.07\% of the gradient. The README puts it as
the targets being ``$<$ 0.1\% of pixels'' that ``1000 steps of a 5.9\,M-parameter model barely see''. Be precise about
``barely see'': it is the \emph{weight} that is small, not the number of visits. $1000 \text{ steps} \times 8 \text{ tiles}
\times 128 \times 128 = 131{,}072{,}000$ pixel-visits, about $131{,}072{,}000 / 1{,}440{,}000 = 91$ per pixel of the scene:
each target is in a tile many times, but each time it is one voice in 16,384.

So what went wrong? Something that affects every pixel. Nudging the network towards this scene's background moved the
residual of every pixel a little, and the careful ranking at the very top got slightly shuffled. No gate on 4.55\% of the
pixels can protect against a change to the other 95\%.

\paragraph{Lesson 2: a pretrained model has no early window.}
GT-HAD's curve - better at 5 epochs than 150 - is a statement about a network that \emph{started from nothing on this scene}.
Early on it has learned the common background and not the targets; that moment is what early stopping finds. Place FT04 on
the same ruler:
\begin{center}
\small
\begin{tabular}{lc}
\toprule
 & detection at $10^{-3}$ false alarms \\
\midrule
GT-HAD, 150 epochs (paper's setting) & 2.0\% \\
GT-HAD, 5 epochs (its best, a diagnostic) & 5.5\% \\
FT04, no adaptation & 6.5\% \\
\bottomrule
\end{tabular}
\end{center}
(GT-HAD figures are PD on valid pixels; FT04's is TPR at 0.1\% with cloud excluded. Close enough to compare, not identical:
different band preparation, and GT-HAD did not exclude cloud. Without exclusion FT04's figure is 6.1\%, still ahead.) FT04
already knows what ordinary land, water and roofs look like, from many scenes, and was never trained on this scene's targets.
It is already where GT-HAD's early stopping is trying to get to, and beyond it. \textbf{A pretrained model starts past the
point early stopping finds} - which is why the damage is the same at $N = 50$ as at $N = 1000$. Any adaptation moves it away
from a good starting point, not towards one.

\begin{keypoint}
Gating and early stopping protect a network that is about to memorise its targets; a pretrained model adapted for a few
hundred steps was never about to, so they have nothing to protect.
\end{keypoint}

\begin{drift}
All results in this chapter are dated 2026-10-19 (scoring gains) and 2026-10-20 (scene adaptation), on FT04. The next
model-side attempts - the uncertainty head (FT06), gated training losses and block masking (FT07, FT08), the progressive
decoder (FT09) and Indradhanu Heavy - are \cref{ch:uncertainty,ch:decoderwork,ch:heavy}.
\end{drift}

\begin{confusions}
\confusion{pAUC went up, so A1 is better.}{pAUC and TPR count pixels. A1 filled in more pixels of objects already found and
lost objects at every size. The analyst reviews objects, not pixels.}
\confusion{The blur in A1 is the same as bench-v2.1's PSF blur.}{Same shape, different job. bench-v2.1 blurs the implant
before scoring to imitate a real target; A1 blurs the residual of the whole scene. That overlap is why A1's bench-v2.1
numbers were treated as partly circular.}
\confusion{A1+ met the criteria, so it should have shipped.}{The criteria were fixed in advance, which is right, but they
checked only pixel measures. A1+ lost 2.7 points of real objects. The record changed the criteria rather than shipping a
regression.}
\confusion{Rank max is the generous choice - a pixel only needs one scorer to like it.}{That generosity is the problem: any
pixel either scorer wrongly likes reaches the top too.}
\confusion{A3 won a benchmark, so it is a better scorer the real bench is too small to show.}{The real bench lost a third of
its pAUC, and the mechanism is understood: real targets cluster in small built-up types and become their types' normal. The
benchmark's placement is what is wrong here.}
\confusion{The mixture scorer is a Gaussian mixture model.}{It is hard k-means types handed to class-aware residual
Mahalanobis: every responsibility is 0 or 1.}
\confusion{The gate made no difference, so gating is useless.}{It made no difference here, because the network did not
memorise the targets in 1000 steps. In a from-scratch method trained for 150 epochs, it does, and gating has a job.}
\confusion{Eval mode means the model is not training.}{\texttt{model.eval()} only freezes BatchNorm's stored statistics and
turns dropout off. The optimiser still updates the weights every step.}
\confusion{Objects caught at the top 1\% went up, so adaptation found more.}{The top 0.1\% went down at the same time. Hits
moved from the very top into the rest of the top 1\%; the analyst looks at the very top first.}
\end{confusions}

\begin{checkyourself}
\item\label{ex:scoringgains:1} Matched filter: $\Cov = \begin{psmallmatrix} 1 & 0.5 \\ 0.5 & 1\end{psmallmatrix}$ and
  $\vect{t} = (1, 1)$. Compute $\vect{w}^\star$ by \eqref{eq:scoringgains:mf} and its SNR by \eqref{eq:scoringgains:mfsnr}.
  Compare with the naive $\vect{w} = \vect{t}$. Explain the result.
\item\label{ex:scoringgains:2} Redo the strip of \cref{we:scoringgains:strip} with the target at pixel 11 raised from 4.8 to
  5.2. Compute its blurred score and say whether it enters the blurred top 4, and what that does to objects caught.
\item\label{ex:scoringgains:3} In \cref{we:scoringgains:rank}, change P1's PCA-RX score from 8 to 11. Recompute both ranks of
  P1 and P4, the rank mean and the rank max. Which pixel does each rule put first?
\item\label{ex:scoringgains:4} In \cref{we:scoringgains:gmm}, compute the responsibilities and the weighted score for a pixel
  at $x = 2$.
\item\label{ex:scoringgains:5} In \cref{we:scoringgains:types}, the built-up type gains two more ordinary roofs at $-2$ and
  $+2$ (8 pixels, still two targets). Recompute the built-up variance and the target score. Do the targets now beat the field
  edges' 3.77?
\item\label{ex:scoringgains:6} Redo the leave-one-scene-out of \cref{we:scoringgains:loso} with only $N = 50$ and $N = 200$
  as candidates. Which $N$ does each scene pick, and what is the headline?
\item\label{ex:scoringgains:7} Using \eqref{eq:scoringgains:dilate}, how many pixels does a gated $1 \times 4$ line cover after
  two passes? What share of a scene would the gate cover if its top 1\% were made entirely of such lines?
\end{checkyourself}
\endgroup

\nextchapter{ch:uncertainty}
